\documentclass[11pt]{article}
\usepackage[utf8]{inputenc}
\usepackage[T1]{fontenc}
\usepackage{lmodern}
\usepackage{geometry}
\usepackage{microtype}
\usepackage{amsmath,amssymb,amsthm,mathtools,bm}
\usepackage{physics}
\usepackage{graphicx}
\usepackage{xcolor}
\usepackage{booktabs}
\usepackage{enumitem}
\usepackage{hyperref}
\usepackage[nameinlink,capitalise]{cleveref}
\geometry{margin=1in}
\hypersetup{colorlinks=true,linkcolor=blue,urlcolor=blue,citecolor=blue}

\newtheorem{definition}{Definition}
\newtheorem{proposition}{Proposition}
\newtheorem{theorem}{Theorem}
\newtheorem{lemma}{Lemma}
\newtheorem{remark}{Remark}

\newcommand{\R}{\mathbb{R}}
\newcommand{\C}{\mathbb{C}}
\newcommand{\Z}{\mathbb{Z}}
\newcommand{\T}{\mathbb{T}} % unit circle
\newcommand{\CP}{\mathbb{C}\mathrm{P}^1}
\newcommand{\RP}{\mathbb{R}\mathrm{P}^1}
\newcommand{\ud}{\mathrm{d}}
\newcommand{\e}{\mathrm{e}}
\newcommand{\ii}{\mathrm{i}}
\newcommand{\E}{\mathbb{E}}
\newcommand{\Var}{\mathrm{Var}}
\newcommand{\Cov}{\mathrm{Cov}}
\newcommand{\sgn}{\mathrm{sgn}}

\title{\textbf{Temporal Coiling--Uncoiling (TCU):\\
A Unified Mixed--Radix--Graph Spectral Framework with Projective Time and Coiling Conservation}}
\author{}
\date{\today}

\begin{document}
\maketitle

\begin{abstract}
We present a self--contained formulation of the \emph{Temporal Coiling--Uncoiling} (TCU) framework for multi--level cascade dynamics.
The theory models discrete level--digits evolving on a mixed--radix lattice, but estimated as smooth circular densities via principled kernels derived from diffusion on the circle and causal time processes.
Time is compactified to $\RP$ to identify $0$ with $\infty$ (``unity through infinity''), making the \emph{unity$\leftrightarrow$spinor} passage a geometric duality.
Three complementary spectral objects extract structure: (i) a radix $L$--function built from band--limited circular moments with a data--driven mirror symmetry; (ii) a modular partition function $Z_{\mathrm{mod}}$ with an operational inverse temperature $\beta_T$ controlling coiling disorder; and (iii) a graph (Ihara--Ruelle) zeta on the finite carry transition graph.
A maximum--entropy coupling of their log--spectra defines the master functional $\Xi$, with cross--cumulants that vanish under independence and quantify multi--view synergy.
We specify all kernels, estimators, and model--selection rules (BIC), and reconcile discrete digits with continuous densities via an explicit back--projection.
Throughout, we emphasize physical motivation, parsimony, and measurable predictions.
Concrete numerical anchors are provided from representative runs (base rate $\lambda_0$, per--level rates, and tail probabilities), and a minimal set of parameters required for reproducible application.
\end{abstract}

\section{Physical motivation and scope}
Many complex systems exhibit bursty, level--resolved cascades with phase--like organization and intermittent \emph{carries} (sudden transitions propagating to higher levels).
Examples include: (i) pattern--forming media with nested oscillators; (ii) switching cascades in mixed--radix encodings of spatially extended time series; and (iii) hierarchical counters in computation or cellular automata.
The TCU framework is designed to (a) infer circular phase densities at each level, (b) quantify carry--conditioned tail risks, and (c) extract spectral invariants across three representations (radix--$L$, modular, and graph--zeta).
It is agnostic to specific microphysics; instead it provides a \emph{data--anchored}, reproducible pipeline with clear hypotheses to falsify.

\paragraph{Operational temperature.}
The inverse temperature $\beta_T$ in $Z_{\mathrm{mod}}(i\beta_T)$ is \emph{operational}: it controls the dispersion of level phases around their means (a measure of coiling disorder), and can be mapped to noise intensity or variability in inter--event statistics (see \cref{sec:modular}).

\paragraph{Parsimony.}
Not every system requires all three spectra. A minimal analysis may use only (a) carry--conditioned tails and the radix $L$--spectrum when transitions are level--centric; or only (b) the graph--zeta when topology dominates; or (c) the modular sector when clear temperature--like control is present. The full $\Xi$ is warranted when cross--covariances between these views are non--negligible.

\section{Geometry of time and the unity$\leftrightarrow$spinor correspondence}
\label{sec:geom}
\subsection{Projective time}
We represent time on the projective line $\RP\simeq S^1/\{\theta\sim \theta+\pi\}$ so that $t=0$ and $t=\infty$ are identified.
Introduce the Cayley map $u\in(-\pi,\pi]\mapsto \tau=\tan(u/2)\in\R\cup\{\infty\}$.
We \emph{choose} the \emph{flat} line element $ds^2=du^2$ on $\RP$; this is equivalent to pulling back the flat metric on the circle by the quotient map, avoiding the $1/\cos^4(u/2)$ factor of stereographic projection. This choice fixes units so that causal kernels are exponentially decaying in $u$, see \cref{sec:kernels}.

\subsection{Coiling field and coiling conservation}
Each level $\ell=0,\dots,L-1$ has a circular phase field $\theta_\ell(x,t)\in\T$ and a smooth density $\rho_\ell(\theta,t|x)\ge 0$ with $\int_{0}^{2\pi}\rho_\ell\,\ud\theta=1$.
Define the \emph{coiling density} $c_\ell(\theta,t|x)=\rho_\ell(\theta,t|x)$ and the \emph{coiling current}
\begin{equation}
  j_\ell(\theta,t|x)= v_\ell(\theta,t|x)\,\rho_\ell(\theta,t|x),\qquad v_\ell(\theta,t|x)=\partial_t \theta_\ell(x,t).
\end{equation}
\begin{definition}[Coiling conservation]
The \emph{coiling balance} at level $\ell$ is
\begin{equation}
  \partial_t \rho_\ell + \partial_\theta j_\ell \;=\; s_\ell^+ - s_\ell^- \;+\; \mathcal{I}_{\ell-1\to \ell} - \mathcal{I}_{\ell\to \ell+1},
  \label{eq:balance}
\end{equation}
where $s^\pm_\ell$ are birth/annihilation source terms (chirality creation in $\pm$ pairs), and $\mathcal{I}_{\ell\to\ell+1}$ denotes \emph{carry flux} (outflow from $\ell$ into $\ell+1$). Summing \eqref{eq:balance} over $\ell$ gives a global conservation law with only boundary terms at the top level (wrap or sink; \cref{sec:wrap}).
\end{definition}

\begin{definition}[Unity and spinor states]
A \emph{unity} state is a maximally coherent configuration with all level phases aligned modulo the radices (a perfectly ``coiled'' reference). A \emph{spinor} state is a minimal stable excitation whose double traversal of $\T$ returns to unity (double cover). The unity$\leftrightarrow$spinor equivalence is the statement that the quotient of phase space by $2\pi$--rotations induces a $\mathbb{Z}_2$ double covering with spinorial monodromy; see \cref{sec:spinor}.
\end{definition}

\section{Levels, digits, and the mixed--radix lattice}
We work with radices $r_\ell\ge 2$ (to be selected by BIC; \cref{sec:BIC}) and discrete digits $s_\ell(x,t)\in\{0,1,\dots,r_\ell-1\}$.
A \emph{carry event} at level $\ell$ is a transition $s_\ell\to s_\ell+1 \pmod{r_\ell}$ accompanied (for carry--exact) by $s_k=0$ for all $k<\ell$ in a guard window. We call a parent at level $\ell$ \emph{eligible} if it satisfies the temporal gap from births and (for carry--exact) the zero--digit guard.
Empirical \emph{parent rates} $\lambda_\ell$ are estimated as counts per observation time; we use weights $w_\ell=\lambda_\ell/\lambda_0$ (dimensionless).

\section{From digits to smooth circular densities}
\label{sec:kernels}
\subsection{Back--projection from discrete digits}
Define disjoint intervals $I_{m}^{(\ell)}=\left[2\pi\frac{m-1/2}{r_\ell},2\pi\frac{m+1/2}{r_\ell}\right)$.
Given empirical digit probabilities $\pi_\ell(m|t,x)$ at time $t$ and site $x$, the \emph{piecewise--constant} density is
\begin{equation}
\rho_\ell^{\mathrm{pc}}(\theta,t|x)=\sum_{m=0}^{r_\ell-1}\frac{\pi_\ell(m|t,x)}{|I_m^{(\ell)}|}\,\mathbf{1}_{\theta\in I_m^{(\ell)}},\qquad |I_m^{(\ell)}|=\frac{2\pi}{r_\ell}.
\end{equation}

\subsection{Smoothing via principled kernels}
We smooth the noisy $\rho_\ell^{\mathrm{pc}}$ by convolution in $(t,x,\theta)$ with the kernels
\begin{align}
K_t(\Delta t;\tau_c)&=\frac{1}{\tau_c}\,\e^{-\Delta t/\tau_c}\,\mathbf{1}_{\Delta t\ge 0},\qquad \tau_c=\frac{c}{\lambda_0},\label{eq:Kt}\\
K_x(\Delta x;\sigma_x)&=\frac{1}{\sqrt{2\pi}\sigma_x}\,\e^{-\Delta x^2/(2\sigma_x^2)},\\
K_\theta(\Delta\theta;\kappa,\phi)&=\frac{1}{2\pi I_0(\kappa)}\,\exp\!\bigl[\kappa\cos(\Delta\theta-\phi)\bigr],\label{eq:Ktheta}
\end{align}
where $I_0$ is the modified Bessel function.
The exponential $K_t$ is the Green function of a causal first--order linear system; the scale $\tau_c$ is fixed by nondimensionalization with speed $c$ (we set $c=1$ units) and base rate $\lambda_0$.
$K_x$ is the heat kernel in one dimension.
$K_\theta$ is the density of a von Mises (VM) distribution on the circle. Its origin is a Fokker--Planck diffusion with potential
\begin{equation}
U_\ell(\theta)=-\kappa_\ell \cos\!\bigl(r_\ell\theta-\phi_\ell\bigr),
\end{equation}
which models $r_\ell$--fold pinning; the stationary solution is VM with concentration $\kappa_\ell$ and mean direction $\phi_\ell/r_\ell$. The \emph{phase bias} $\phi_\ell$ is estimated from the first circular moment $M_{1,\ell}=\int \e^{\ii\theta}\rho_\ell\,\ud\theta$ via $\phi_\ell=\arg(M_{1,\ell}^{r_\ell})$.
The concentration $\kappa_\ell$ is inferred from $|M_{1,\ell}|$ using the standard VM moment relation $|M_{1,\ell}|=I_1(\kappa_\ell)/I_0(\kappa_\ell)$.

\subsection{Band--limited circular moments}
Define circular Fourier coefficients
\begin{equation}
\widehat{\rho}_\ell(k;t,x)=\int_{0}^{2\pi}\rho_\ell(\theta,t|x)\,\e^{\ii k\theta}\,\ud\theta,\qquad k\in\Z.
\end{equation}
We adopt a \emph{bandlimit} $K_\ell$ chosen by BIC (below) and retain only $|k|\le K_\ell$.
For discrete digits the natural basis is the finite set $k=0,1,\dots,r_\ell-1$, since higher modes alias under the partition $I_m^{(\ell)}$.

\section{Spectral object I: the radix $L$--function}
\label{sec:Lradix}
For each level $\ell$ and retained mode $k$, define the normalized coefficient
\begin{equation}
\alpha_{\ell,k}(t)=\frac{\widehat{\rho}_\ell(k;t)}{\widehat{\rho}_\ell(0;t)}\in\C,\qquad \alpha_{\ell,0}=1.
\end{equation}
The radix $L$--function is
\begin{equation}
L_{\mathrm{radix}}(s;t)=\prod_{\ell=0}^{L-1}\prod_{k\in \mathcal{K}_\ell}\bigl(1-\alpha_{\ell,k}(t)\,r_\ell^{-s}\bigr)^{-w_\ell},\qquad \mathcal{K}_\ell=\{1,\dots,K_\ell\},
\label{eq:Lradix}
\end{equation}
with weights $w_\ell=\lambda_\ell/\lambda_0$.
\paragraph{Mirror symmetry and palindromy.}
If the radix list is palindromic up to weights (i.e.\ $r_\ell=r_{L-1-\ell}$ and $w_\ell=w_{L-1-\ell}$), and if the first moments obey a real phase symmetry (empirically, $\alpha_{\ell,k}\approx \overline{\alpha}_{L-1-\ell,k}$), then $L_{\mathrm{radix}}$ admits an \emph{approximate functional equation}
\begin{equation}
\Lambda(s;t):=Q(s)\,L_{\mathrm{radix}}(s;t)\;\approx\;\varepsilon(t)\,\Lambda(1-s;t),
\end{equation}
for an explicit normalizing factor $Q(s)$ and a data--estimated $\varepsilon(t)\in\C$ with $|\varepsilon|\le 1$ (from the argument of $\sum_m \pi_\ell(m)\,\e^{2\pi\ii m/r_\ell}$).
Mirror is \emph{tested}, not imposed; if rejected, we simply analyze $L_{\mathrm{radix}}$ on the critical line $\Re s=1/2$ without symmetry claims.

\section{Spectral object II: modular partition function}
\label{sec:modular}
Let $\beta_T>0$ be an operational inverse temperature controlling phase dispersion.
Associate to each level an effective stiffness $\kappa_\ell$ from \cref{sec:kernels} and define level spectra $E_{\ell,n}=\frac{n^2}{2\kappa_\ell}$, $n\in\Z$ (harmonic approximation of small phase fluctuations).
The modular partition is then
\begin{equation}
Z_{\mathrm{mod}}(\tau)=\prod_{\ell=0}^{L-1}\vartheta_3\!\left(0\,\middle|\,\frac{\ii}{2\pi}\beta_T \kappa_\ell\right)^{w_\ell}\,\eta(\ii \beta_T)^{-\gamma},
\qquad \tau=\ii\beta_T,
\end{equation}
with Jacobi $\vartheta_3$ and Dedekind $\eta$, and a small integer $\gamma$ collecting residual degrees of freedom.
Under the modular $S$--transform $\tau\mapsto -1/\tau$ one has the duality
\begin{equation}
Z_{\mathrm{mod}}(-1/\tau)\;=\;\prod_\ell \bigl(\sqrt{\kappa_\ell/\beta_T}\bigr)^{w_\ell}\,Z_{\mathrm{mod}}(\tau)\;\times\; (\text{phase}).
\end{equation}
We interpret a \emph{self--dual point} as a balance between coiling (order) and uncoiling (disorder).

\section{Spectral object III: finite--graph (Ihara--Ruelle) zeta}
\label{sec:ihara}
Build a directed multigraph $G=(V,E)$ whose vertices are level--digit states and whose edges are observed carry transitions (optionally weighted by rates and phases).
Let $A$ be the (weighted) adjacency and $Q$ the diagonal with $Q_{vv}=d_v-1$ (outdegree minus one).
The Ihara zeta is
\begin{equation}
\zeta_G(u)=\prod_{[P]}\bigl(1-u^{\ell(P)}\bigr)^{-1},
\end{equation}
product over \emph{prime cycles} (non--backtracking cycles, counted up to cyclic rotation).
It admits the determinant formula (for finite $G$)
\begin{equation}
\zeta_G(u)^{-1}=(1-u^2)^{|E|-|V|}\det\!\bigl[I - u A + u^2 Q\bigr].
\end{equation}
We use $u=\e^{-z}$ so that $\Re z>0$ corresponds to absolutely convergent series, and we allow a \emph{twist} by complex edge phases from level moments, producing a \emph{twisted} zeta. Poles quantify recurrent carry patterns and their stability.

\section{Master functional $\Xi$ from maximum--entropy coupling}
Let
\begin{equation}
X_1=\log|L_{\mathrm{radix}}(s)|,\qquad X_2=\log|Z_{\mathrm{mod}}(\tau)|,\qquad X_3=\log|\zeta_G(\e^{-z})|.
\end{equation}
Define their centered versions $\tilde X_i=X_i-\E X_i$ over a sampling ensemble (sites/time windows).
Assuming (i) small, stationary log--fluctuations and (ii) no further structure beyond measured second moments, the maximum--entropy density over $(\tilde X_1,\tilde X_2,\tilde X_3)$ is Gaussian with covariance $\Sigma$.
This yields the \emph{quadratic} generating functional
\begin{equation}
\Xi(s,\tau,z)=\exp\!\left\{\sum_{i}\alpha_i \tilde X_i+\frac{1}{2}\sum_{i,j} C_{ij}\,\tilde X_i\,\tilde X_j\right\},
\label{eq:Xi}
\end{equation}
with $\alpha_i$ and $C_{ij}$ determined empirically (e.g.\ by method of moments). The cross--terms $C_{ij}$ precisely measure departure from independence across the three views.

\section{Model selection, bandlimits, and palindromy test}
\label{sec:BIC}
Let $\mathcal{M}$ index candidate radix lists $\bm r=(r_0,\dots,r_{L-1})$, bandlimits $\bm K=(K_0,\dots,K_{L-1})$, and wrap/sink choice at the top.
Given data $\mathcal{D}$ (eligible parents per level; \emph{eligible} means beyond the birth exclusion window and satisfying carry--exact guard for the conditioned analysis), define the log--likelihood of the smoothed densities and carries under the kernels of \cref{sec:kernels}.
Use the Bayesian information criterion
\begin{equation}
\mathrm{BIC}(\mathcal{M})=-2\log L_{\max}(\mathcal{D}|\mathcal{M})+k(\mathcal{M})\log N(\mathcal{D}),
\end{equation}
with $N(\mathcal{D})=\sum_\ell n_\ell$ the total number of eligible parents and $k(\mathcal{M})$ the number of free parameters (including bandlimits and kernel scales).
A palindromy \emph{test} compares the best palindromic model to the best general model (likelihood ratio with BIC penalty); if rejected, no mirror is used in \cref{sec:Lradix}.

\section{Estimation and empirical anchors}
\paragraph{Weights from rates.}
Use $w_\ell=\lambda_\ell/\lambda_0$, where $\lambda_\ell$ are per--level parent rates estimated from counts.
For a representative dataset we observed a pooled base rate $\lambda_0\approx 0.475$ over total time $T\approx 1.28\times 10^4$, with per--level rates $\lambda_0\approx 0.475$, $\lambda_1\approx 0.158$, $\lambda_2\approx 0.153$, $\lambda_3\approx 0.117$, and tail probabilities $p^{(2)}_\ell\in\{0.333,0.965,0.763\}$ at lower levels (numbers rounded).%
\footnote{These anchors reflect an exemplar run summarized in \texttt{rates\_all.json}.}
A typical coil--$L$ scan used $L=4$, geometric weights, $10$ evenly spaced sample times in $[0,200]$, and $2\cdot 10^4$ sampled pairs per time.%
\footnote{Representative scan configuration summarized in \texttt{weights.json}.}

\paragraph{Carry--conditioned tails and $\beta_\ell$ slopes.}
For each level estimate a Jeffreys--smoothed tail probability conditioned on parent eligibility. Regress $\log p^{(2)}_\ell$ on $\ell$ to obtain slopes $\beta_\ell$ (distinct from $\beta_T$). These slopes calibrate the radius of convergence for the product in \eqref{eq:Lradix} via the mapping $K_\ell\leftrightarrow$ effective bandwidth.

\section{Unity--spinor equivalence in the projective picture}
\label{sec:spinor}
Let $\pi:\T\to \RP$ be the double cover $[\theta]\mapsto [\theta]\ (\mathrm{mod}\ \pi)$. The fundamental group of $\RP$ is $\Z_2$, and lifting a closed loop in projective time to $\T$ produces a sign change after one circuit (spinorial behavior). We say a configuration is \emph{spinor--like} if its lifted phase requires a $4\pi$ rotation to return to unity; otherwise it is \emph{unity--like}. This formalizes unity$\leftrightarrow$spinor as a geometric dual, not a dynamical claim.

\section{Top--level boundary: wrap vs.\ sink}
\label{sec:wrap}
In \eqref{eq:balance} the top carry flux $\mathcal{I}_{L-1\to L}$ is closed either by \emph{wrap} (identifying level $L$ with level $0$ modulo a winding counter $\mathsf{w}\in\Z$ that records net uncoiling) or by a \emph{sink} (absorbing boundary, modeling loss). The wrap realizes the $0\equiv\infty$ identification on $\RP$ at the level hierarchy.

\section{When is the full $\Xi$ needed?}
If the covariance matrix $\Sigma$ of $(\tilde X_1,\tilde X_2,\tilde X_3)$ is close to diagonal, individual spectra suffice (parsimony). Significant off--diagonal entries justify the master coupling \eqref{eq:Xi}. In practice, a held--out predictive score (e.g.\ cross--validated log score) can decide whether the extra parameters $C_{ij}$ are warranted.



\section{Methods: event selection, coding, and diagnostics}

\paragraph{Phase and mixed–radix coding.}
We construct the $Z$–candidate azimuth $\phi_Z$, map it to $\mathrm{Uniform}[0,1)$ via an empirical CDF built from the same sample, and expand it in mixed radices $[3,5,17,5,3]$ to obtain digits $s_0,\dots,s_4$. We also bin the average lepton azimuth into 8 equal “site’’ bins. Quality–control thresholds are: $\chi^2$ uniformity at $\alpha=0.01$, pairwise MI threshold $0.02$\,bits, and palindromy JSD threshold $0.02$\,bits. The $\mu\mu$ flattening summary records a total weight of $4{,}964{,}769$, \texttt{emit\_u=true}, and the path to the emitted uniformized time series used in the diagnostics. :contentReference[oaicite:7]{index=7}

\paragraph{Exporter (ROOT$\to$CSV) and reproducible commands.}
We export \texttt{triad\_timeseries.csv} and \texttt{pairs.csv} with \texttt{runNumber,eventNumber} metadata (and \texttt{lumiBlock} only if present across all inputs).
The following commands reproduce the results:

\noindent\textbf{Electrons ($ee$):}
\begin{verbatim}
python atlas_to_tcu_v5.py \
  --input data_A.2lep.root data_B.2lep.root data_C.2lep.root data_D.2lep.root \
  --out out_phiZ_v5 --sample ee --require-trigger \
  --phase-source phiZ --write-pairs --emit-theta

python pairs_to_uniform_digits_v3.py \
  --pairs out_phiZ_v5/pairs.csv \
  --column phiZ --weights none \
  --radices 3 5 17 5 3 --n-sites 8 \
  --out out_uniform_phiZ_v5 --emit-u

python tcu_diagnostics.py \
  --csv out_uniform_phiZ_v5/triad_timeseries_uniform.csv \
  --radices 3 5 17 5 3 \
  --pairs out_phiZ_v5/pairs.csv \
  --report out_uniform_phiZ_v5/tcu_report.json \
  --alpha 0.01 --mi-threshold 0.02 --jsd-threshold 0.02 \
  --bandlimit 0.10 --band-frac 0.70
\end{verbatim}
The resulting \texttt{tcu\_report.json} records $n=4{,}418{,}648$, digit uniformity (all $p$ high), pairwise MI $\le 2.3\times 10^{-6}$\,bits, JSD $=8.81\times 10^{-5}$\,bits, and a $10\%$ low–band fraction $0.09965$. :contentReference[oaicite:8]{index=8}
A per–run bandlimit pass (sorted by \texttt{eventNumber} within each run; $10\%$ band) gives the weighted overall $0.099716$. :contentReference[oaicite:9]{index=9}
For a pre–flattening control on the same selection, the most–significant digit fails uniformity with $p\simeq 1.6\times 10^{-25}$, even though MI and palindromy are already small. :contentReference[oaicite:10]{index=10}

\noindent\textbf{Muons ($\mu\mu$):}
\begin{verbatim}
python atlas_to_tcu_v5.py \
  --input data_A.2lep.root data_B.2lep.root data_C.2lep.root data_D.2lep.root \
  --out out_muZ_v5 --sample mumu --require-trigger \
  --phase-source phiZ --write-pairs --emit-theta

python pairs_to_uniform_digits_v3.py \
  --pairs out_muZ_v5/pairs.csv \
  --column phiZ --weights none \
  --radices 3 5 17 5 3 --n-sites 8 \
  --out out_uniform_muZ_v5 --emit-u

python tcu_diagnostics.py \
  --csv out_uniform_muZ_v5/triad_timeseries_uniform.csv \
  --radices 3 5 17 5 3 \
  --pairs out_muZ_v5/pairs.csv \
  --report out_uniform_muZ_v5/tcu_report.json \
  --alpha 0.01 --mi-threshold 0.02 --jsd-threshold 0.02 \
  --bandlimit 0.10 --band-frac 0.70
\end{verbatim}
The \texttt{flatten\_summary.json} (for the $\mu\mu$ run) records the CDF settings and total weight; the \texttt{tcu\_report.json} shows $n=4{,}964{,}769$, digit uniformity (all $p$ high), pairwise MI $\le 2.5\times 10^{-6}$\,bits, JSD $=8.42\times 10^{-5}$\,bits, and a $10\%$ low–band fraction $0.10048$. 
A per–run bandlimit pass (sorted by \texttt{eventNumber} within each run) gives the weighted overall $0.100394$ in a $10\%$ band. :contentReference[oaicite:12]{index=12}

\paragraph{Temporal checks.}
Sorting globally by \texttt{(run,event)} leaves the $10\%$ low–band fraction near the band size in both channels (verbatim console: \texttt{ee: 0.099763}, \texttt{$\mu\mu$: 0.100397}). In the absence of a unified detector–time coordinate across all files, we report these temporal tests as inconclusive rather than falsifying.



\section{Results: ATLAS $Z\!\to\!\ell\ell$ digit tests}

We test the TCU mixed–radix code on ATLAS $Z\!\to\!\ell\ell$ in both $ee$ and $\mu\mu$ channels, using the $Z$-candidate azimuth $\phi_Z$ mapped to a uniform circle by an empirical CDF, and the radices $[3,5,17,5,3]$. On multi–million–event samples, the \emph{distributional} predictions (digit uniformity, pairwise independence across digit positions, and palindromy) are decisively confirmed; the \emph{temporal} (bandlimit) prediction is inconclusive under the available orderings (unsorted stream and global sort by \texttt{(run,event)}). 

\paragraph{Electrons ($Z\!\to\!e^+e^-$).}
With $n=4{,}418{,}648$ events, all five digits pass $\chi^2$ uniformity (three digits have $p\simeq 1$; the others $p=0.983,0.905$), pairwise mutual information across positions is $\le 2.3\times 10^{-6}$\,bits, and the palindromy divergence is $8.81\times 10^{-5}$\,bits. The low–frequency energy in a $10\%$ band (unsorted stream) is $0.09965\approx$\,band size. These values are taken from the \texttt{tcu\_report.json} produced by the pipeline. :contentReference[oaicite:0]{index=0}

A per–run analysis (sorting by \texttt{eventNumber} within each run, $10\%$ band, minimum 10k events per run) yields a weighted overall of $0.099716$, again consistent with a white baseline at the chosen band size. :contentReference[oaicite:1]{index=1}

For comparison, a pre–flattening control on the same sample shows the expected MSD bias: the most–significant digit fails uniformity with $\chi^2=114.15$ (dof=2), $p\simeq 1.6\times 10^{-25}$, while MI and palindromy are already small—demonstrating the necessity of empirical flattening of $\phi$. :contentReference[oaicite:2]{index=2}

\paragraph{Muons ($Z\!\to\!\mu^+\mu^-$).}
With $n=4{,}964{,}769$ events, all five digits pass $\chi^2$ uniformity (three digits have $p\simeq 1$; the others $p=0.912,0.699$), pairwise mutual information across positions is $\le 2.5\times 10^{-6}$\,bits, and the palindromy divergence is $8.42\times 10^{-5}$\,bits. The $10\%$ low–frequency band fraction (unsorted stream) is $0.10048\approx$\,band size. (See \texttt{tcu\_report.json}.) :contentReference[oaicite:3]{index=3}
The flattening summary records a total weight of $4{,}964{,}769$, \texttt{emit\_u=true}, and the emitted uniformized time series path used below. :contentReference[oaicite:4]{index=4}
The per–run bandlimit (sorted by \texttt{eventNumber} within each run) gives a weighted overall of $0.100394$ in a $10\%$ band. :contentReference[oaicite:5]{index=5}

\paragraph{Temporal ordering checks.}
Sorting globally by \texttt{(run,event)} leaves the $10\%$ low–band fraction near the band size in both channels (verbatim console): 
\texttt{ee: 0.099763 (band=0.1)}, \texttt{$\mu\mu$: 0.100397 (band=0.1)}.
A band–size scan on $\mu\mu$ is near–diagonal (\,$0.02 \!\to\! 0.020083$, $0.05 \!\to\! 0.050196$, $0.10 \!\to\! 0.100397$, $0.20 \!\to\! 0.200147$, $0.30 \!\to\! 0.299817$\,), consistent with a white baseline under the available ordering.
Overall, the TCU \emph{distributional} predictions are confirmed, while the temporal prediction is \emph{inconclusive} here due to the absence of a detector–time coordinate common to all files. 

\begin{table}[t]
\centering
\caption{Summary of key metrics for $ee$ and $\mu\mu$ (band$=10\%$). ``Per–run overall'' is the weighted mean after sorting by \texttt{eventNumber} within runs.}
\label{tab:atlas-summary}
\begin{tabular}{lcccccc}
\toprule
Channel & $n$ & Max MI (bits) & Palindr.\ JSD (bits) & Low–band (unsorted) & Per–run overall \\
\midrule
$ee$   & 4{,}418{,}648 & $\le 2.3\times 10^{-6}$ & $8.81\times 10^{-5}$ & 0.09965 & 0.099716 \\
$\mu\mu$ & 4{,}964{,}769 & $\le 2.5\times 10^{-6}$ & $8.42\times 10^{-5}$ & 0.10048 & 0.100394 \\
\bottomrule
\end{tabular}
\end{table}



\section{Summary of assumptions and how to falsify them}
\begin{itemize}[leftmargin=2em]
\item \textbf{Projective time:} $ds^2=du^2$ on $\RP$ (Cayley chart). Falsify by showing causal kernels demand non--exponential tails in $u$.
\item \textbf{VM smoothing:} arises from $r_\ell$--fold pinning potential; falsify if empirical circular moments are inconsistent with a VM family.
\item \textbf{Bandlimit:} $K_\ell$ finite and BIC--selected; falsify if higher harmonics persist under penalized likelihood.
\item \textbf{Mirror:} tested, not assumed; reject if $\varepsilon$ significantly complex and asymmetric.
\item \textbf{Ihara zeta:} carry graph finite and non--backtracking cycles well--defined; falsify if transitions require higher--order memory (then lift the graph).
\end{itemize}

\appendix
\section{Derivation of the VM kernel from Fokker--Planck}
Consider $\partial_t \rho=\partial_\theta\bigl(D\partial_\theta \rho + \rho\,U'(\theta)\bigr)$ with $U(\theta)=-\kappa\cos(r\theta-\phi)$.
A stationary solution satisfies $D\partial_\theta \rho + \rho\,U'(\theta)=0$, hence
\begin{equation}
\frac{\partial_\theta \rho}{\rho}=-\frac{U'(\theta)}{D}\;\Rightarrow\;\rho(\theta)\propto \exp\!\Bigl(\frac{\kappa}{D}\cos(r\theta-\phi)\Bigr),
\end{equation}
which is VM with concentration $\kappa/D$. Absorb $D$ into $\kappa$ by scale choice.

\section{Ihara zeta determinant formula (sketch)}
For a finite undirected graph the Bass--Ihara formula gives
\begin{equation}
\zeta_G(u)^{-1}=(1-u^2)^{|E|-|V|}\det\!\bigl[I - u A + u^2 Q\bigr].
\end{equation}
For directed multigraphs and complex weights, an analogous non--backtracking matrix yields the same structure; see standard references on Ihara zeta generalizations. Twists by $U(1)$ phases enter as unitary similarities.

\section{Discrete--continuous reconciliation}
Given $\pi_\ell(m)$ and $\rho_\ell$ we have the exact identities
\begin{equation}
\pi_\ell(m)=\int_{I_m^{(\ell)}}\rho_\ell(\theta)\,\ud\theta,\qquad
\widehat{\rho}_\ell(k)=\sum_{m=0}^{r_\ell-1}\pi_\ell(m)\,\widehat{\mathbf{1}}_{I_m^{(\ell)}}(k).
\end{equation}
Thus cyclotomic factors in \eqref{eq:Lradix} are justified at finite $k$ by the finite partition.

\section*{Data anchors and reproducibility}
\begin{itemize}
\item Example base rate and per--level rates, with carry--tail summaries, drawn from a representative run (\texttt{rates\_all.json}).
\item A representative coil--$L$ scan configuration (\texttt{weights.json}) records level count, weighting scheme, and sampling schedule.
\end{itemize}
We recommend publishing these JSONs alongside figures for full reproducibility.

\end{document}
